% ============================================================
% PerceptFence — Springer Nature Cybersecurity submission
% Template: sn-jnl.cls with sn-mathphys-num bibliography style
% De-anonymized: Neeraj Kumar Singh, Parafin, Inc.
% Venue: https://www.editorialmanager.com/cyse/default.aspx
% Created: 2026-05-03 (NEE-245)
% QA gate: PASS (NEE-243, 2026-05-02)
% ============================================================

\documentclass[sn-mathphys-num]{sn-jnl}

\usepackage{graphicx}
\usepackage{multirow}
\usepackage{amsmath,amssymb}
\usepackage{booktabs}
\usepackage{array}
\usepackage{url}

%%% ── Title ─────────────────────────────────────────────────────────────────

\title[Consent-Aware Runtime Mediation]{%
  Consent-Aware Runtime Mediation for Privacy in Real-Time
  Screen-Share AI Assistants}

%%% ── Authors and affiliations ──────────────────────────────────────────────

\author*[1]{\fnm{Neeraj Kumar} \sur{Singh}}

\affil*[1]{%
  \orgname{Parafin, Inc.},%
  \orgaddress{%
    \city{San Francisco},%
    \state{California},%
    \country{USA}%
  }%
}

\corresp{*Corresponding author. E-mail: \url{b.neerajkumarsingh@gmail.com}}

%%% ── Abstract ───────────────────────────────────────────────────────────────

\abstract{%
Live screen-share AI assistants can observe raw screen and speech streams,
but users lack fine-grained runtime control over what the assistant may see,
retain, and say.
Static chatbot and privacy controls do not close this gap: the sensitive
content is in the live capture, not in the prompt, and the assistant needs
fast-changing interface context to respond at all.
We design and evaluate a runtime mediation layer that enforces consent,
redaction, memory controls, and output safeguards for live multimodal
assistance, sitting between capture adapters, session memory, and assistant
responses.
We evaluate the layer on synthetic screen-share scenarios covering secrets,
personal data, sensitive speech fragments, screen-visible prompt injection,
dense or zoomed interfaces, fast window switching, and three
adversarial-evasion classes (Unicode-homoglyph credentials, split or spaced
personal-data digits, and role-play\,/\,chain-prompt screen instructions),
measuring sensitive exposure, false blocks, latency overhead, and task
completion against an unmediated prototype baseline.
We contribute a design and evaluation of runtime control mechanisms for
screen-share assistants, with claims bounded to synthetic evidence.}

%%% ── Keywords ───────────────────────────────────────────────────────────────

\keywords{privacy, screen-share assistants, runtime mediation, consent,
          redaction, multimodal AI, prompt injection, evaluation}

\begin{document}

\maketitle

%%%──────────────────────────────────────────────────────────────────────
\section{Introduction}\label{sec:intro}
%%%──────────────────────────────────────────────────────────────────────

Multimodal AI assistants are increasingly deployed inside live screen-share
sessions, where they observe the same surface a human is explaining:
terminals with credentials, browsers with personal data, chat overlays with
personal notifications, and spoken fragments containing sensitive information.
This deployment pattern is fundamentally different from the chatbot setting:
the sensitive content is not in the user's prompt; it is in the live capture
stream itself.

Static privacy settings---blanket ``do not log,'' ``do not share with
vendors'' toggles---address a different problem.
They cannot answer questions like \textit{``this terminal window is okay but
the password manager is not,''} \textit{``summarise what is on screen but
never name the customer,''} or \textit{``redact this email address from your
reply but use it to find the right ticket.''}
These questions are about \textbf{runtime} mediation between capture and
assistant: deciding, frame-by-frame and utterance-by-utterance, what the
model sees, what gets retained, and what comes out.

This paper presents a runtime mediation layer for live screen-share AI
assistants.
Our frozen thesis:

\begin{itemize}
  \item \textbf{Problem.}
        Live screen-share AI assistants can observe raw screen and speech
        streams, but users lack fine-grained runtime control over what the
        assistant may see, retain, and say.
  \item \textbf{What we solve.}
        A runtime mediation layer that enforces consent, redaction, memory
        controls, and output safeguards for live multimodal assistance.
  \item \textbf{Contribution.}
        Design and evaluation of a consent-aware runtime layer for real-time
        screen-share assistants.
\end{itemize}

The layer interposes seven cooperating modules between raw capture and
assistant response---screen capture adapter, speech event adapter, redaction
engine, consent\,/\,policy engine, session memory gate, output guard, and
audit logger---with three control surfaces: \textbf{observe}, \textbf{retain},
and \textbf{say}.
We design the layer to satisfy a stated adversary model
(Section~\ref{sec:threat}), describe per-module enforcement and trust
boundaries (Section~\ref{sec:design}), implement and test the prototype
(Section~\ref{sec:impl}), and evaluate per-module contributions on a synthetic
benchmark of eleven scenario classes including three adversarial-evasion
classes (Section~\ref{sec:eval}).

Our contributions:

\begin{itemize}
  \item \textbf{A consent-aware runtime mediation design} for screen-share AI
        assistants, decomposed into seven modules with explicit per-module
        trust assumptions and enforces-vs.-does-not-enforce boundaries.
  \item \textbf{A synthetic evaluation benchmark} of eleven scenario classes
        covering credentials, personal data, sensitive speech, prompt injection
        on-screen, dense\,/\,zoomed interfaces, fast window switching, and
        three adversarial-evasion classes (Unicode-homoglyph credentials,
        split\,/\,spaced personal-data digits, role-play\,/\,chain-prompt
        screen instructions).
  \item \textbf{A per-module ablation study} isolating which mediation module
        removes or blocks each annotated risk class, including a measurement
        that the full guard achieves the expected outcome on every fixture
        (11\,/\,11) while the unmediated baseline reaches it on none (0\,/\,11).
  \item \textbf{A claim-to-metric audit framework} that ties every empirical
        claim in the evaluation to an annotated metric definition, preventing
        unmeasured ``privacy'' claims that have weakened prior work.
\end{itemize}

We make no claims about formal privacy preservation, differential privacy
guarantees, defence against compromised hosts or models, or completeness
against unknown adversarial-evasion classes.
The contribution is bounded: a design and evaluation of runtime control
mechanisms on a synthetic benchmark.

%%%──────────────────────────────────────────────────────────────────────
\section{Threat Model}\label{sec:threat}
%%%──────────────────────────────────────────────────────────────────────

We separate in-scope runtime threats from excluded systems-security threats.
The contribution is meaningful only relative to the risks it chooses to
handle.

\subsection{Adversary classes}\label{sec:adversaries}

Table~\ref{tab:adversaries} lists the eight adversary classes considered.

\begin{table}[ht]
\caption{Adversary classes and scope decisions.}
\label{tab:adversaries}
\begin{tabular}{@{}lp{3.8cm}lp{5.2cm}@{}}
\toprule
ID & Adversary & In scope? & Rationale \\
\midrule
A1 & Malicious screen content (attacker controls displayed text\,/\,visuals,
     including instructions meant to manipulate the assistant)
   & Yes --- primary
   & Prompt injection through displayed content is the central
     screen-share-specific risk. \\[4pt]
A2 & Malicious meeting participant (deliberately displays sensitive content
     or prompts the assistant to reveal content outside another
     participant's consent)
   & Yes
   & Tests whether consent and output policies hold under adversarial
     participation. \\[4pt]
A3 & Curious or over-privileged user (tries to lower
     consent\,/\,redaction\,/\,retention to expose more than the policy allows)
   & Yes
   & Distinguishes enforceable runtime policy from advisory user
     preferences. \\[4pt]
A4 & Assistant output leakage (model repeats, summarises, infers, or
     indirectly references content that should have been redacted or gated)
   & Yes
   & Motivates the memory gate and output guard. \\[4pt]
A5 & Temporal exposure (sensitive content appears briefly during window
     switches, popups, zoom changes, or small-font rendering)
   & Yes
   & Matches the synthetic scenarios and tests mediation under changing
     context. \\[4pt]
A6 & Compromised infrastructure (read\,/\,modify audit logs, session
     storage, runtime state, backend services)
   & \textbf{No}
   & Reasonable exclusion for the current prototype.
     We do not claim tamper-evidence or infrastructure hardening. \\[4pt]
A7 & Poisoned model or supply chain
   & \textbf{No}
   & Broader systems-security problem named as out of scope. \\[4pt]
A8 & Operating-system compromise (malware controls host OS, display server,
     microphone stack, or clipboard)
   & \textbf{No}
   & Cannot be defended against without a trusted-computing design. \\
\bottomrule
\end{tabular}
\end{table}

\subsection{Trust boundaries}\label{sec:trust}

Captured frames, extracted text, audio events, and notification events are
treated as \textbf{untrusted input} until the policy and redaction stages
finish.
The user's consent profile is \textbf{trusted at session start} and
modifiable only via authenticated policy operations within the session.
Backend storage and audit logs are \textbf{trusted under the A6 exclusion};
we do not provide tamper-evidence beyond an append-only hash chain
(Section~\ref{sec:audit}).

\subsection{Claim guardrails}\label{sec:guardrails}

Every empirical claim in this paper maps to a metric defined in
Section~\ref{sec:metrics} and computed on the benchmark in
Section~\ref{sec:ablation}.
The phrase ``for privacy'' in the title refers to a measured \textit{reduction}
in sensitive exposure rate on synthetic fixtures---not formal privacy
preservation, differential privacy, side-channel absence, or prevention of all
leakage.

%%%──────────────────────────────────────────────────────────────────────
\section{Related Work}\label{sec:related}
%%%──────────────────────────────────────────────────────────────────────

\paragraph{Runtime capture permissions.}
Mobile operating systems use API-level permissions to control which
applications may access the camera, microphone, or screen-capture surface.
Felt et al.\ characterise Android's permission model and find that most users
grant permissions without understanding their scope~\cite{felt2012permissions}.
A companion study shows that applications routinely request permissions they do
not use, creating unnecessary exposure~\cite{felt2011android}.
At the system layer, TaintDroid tracks the flow of privacy-sensitive data
through the Android runtime, demonstrating that third-party applications
frequently transmit sensitive user data to advertising networks without user
awareness~\cite{enck2014taintdroid}.
These system-level designs gate access to a capture API but do not control
what happens to the content of the captured stream after access is granted.
PerceptFence operates at the content layer rather than the access layer: it
interposes between a granted capture stream and an AI assistant's
context-construction step.

\paragraph{Privacy norms and contextual integrity.}
Nissenbaum frames appropriate information flow in terms of contextual
integrity: information should flow in ways consistent with the norms of the
context in which it was originally disclosed~\cite{nissenbaum2004privacy}.
Apthorpe et al.\ apply contextual integrity to smart-home IoT devices, showing
that users hold nuanced norms about which data a home device may forward to
third-party services, conditioned on the recipient and stated
purpose~\cite{apthorpe2018smart}.
Shvartzshnaider and Duddu extend this framework to large language models,
measuring deviations between the information-flow choices made by LLMs and
contextually appropriate expectations, and proposing an auditing metric
grounded in contextual integrity theory~\cite{shvartzshnaider2026privacy}.
The consent and policy engine in PerceptFence enforces a simplified contextual
boundary by content type: it blocks sensitive screen and speech content from
entering the assistant's context unless the active session policy permits that
content category.

\paragraph{Privacy-aware AI assistants.}
Xu et al.\ study user and expert expectations of AI-powered privacy assistants
that automate privacy decisions on behalf of users, identifying factors that
influence acceptability, including transparency of information sources and
regulatory context~\cite{xu2025acceptability}.
Danry et al.\ build a gaze-aware multimodal assistant that uses egocentric video
to infer where a user is struggling during a task, raising questions about what
such an assistant is permitted to observe and how that observation boundary
should be defined~\cite{danry2026gaze}.
These systems address AI assistant design and user-facing controls but do not
impose runtime content-level filtering on a continuously changing multimodal
stream.
PerceptFence contributes a mediation layer targeting live screen-share input
rather than post-hoc user control or assistant behaviour calibration.

\paragraph{Screen-capture privacy in deployed systems.}
Microsoft Recall captures a desktop screenshot every few seconds and uses
on-device LLMs to provide a retrievable memory layer; early deployments did not
filter sensitive content from the snapshot store, and subsequent versions added
opt-in filtering for recognised credential patterns after significant public and
security-research scrutiny~\cite{microsoft2024recall}.
Zoom AI Companion provides in-meeting AI assistance including transcription and
summaries; its privacy documentation describes zero-data-retention options and
per-account toggle controls, but these operate at the session level rather than
at the granularity of what the assistant may observe within an active
session~\cite{zoom2024ai}.
Neither system exposes per-content-category runtime controls during an active
session.
PerceptFence targets this gap: it mediates at the frame and utterance level
rather than only at session setup time.

\paragraph{Prompt injection attacks.}
Screen-visible content introduces a direct prompt injection surface when an AI
assistant observes and processes text displayed on screen.
Greshake et al.\ systematise indirect prompt injection attacks in
LLM-integrated applications, showing that attacker-controlled content placed in
retrieved data can override application instructions without any direct user
interaction~\cite{greshake2023indirect}.
Liu et al.\ extend this analysis to black-box attacks on commercial
applications, reporting high success rates against production LLM integrations
and proposing a taxonomy of injection techniques~\cite{liu2023prompt}.
Screen-share sessions are a particularly direct injection surface because
attacker-controlled text is visually displayed and the assistant observes the
display without intermediation; the synthetic benchmark includes
prompt-injection-on-screen and encoded-screen-instruction scenarios covering
this threat.

\paragraph{Defences against prompt injection and agent safeguards.}
StruQ separates instructions and data into structured channels, fine-tuning the
LLM to follow only the designated instruction channel, reducing injection
success rates substantially on the evaluated benchmark~\cite{chen2025struq}.
SecAlign uses preference optimisation to train an LLM to rank policy-aligned
outputs above injection-compliant alternatives, achieving injection success
rates below ten percent against a range of attacks including ones not seen
during training~\cite{chen2025secalign}.
Yang et al.\ demonstrate that injected content written into an agent's
long-term memory can persist across sessions and trigger unauthorised tool
actions in future sessions~\cite{yang2026zombie}; this result motivates
PerceptFence's memory gate, which prevents sensitive content from being
written to session memory under configurable policy.
Chen and Cong propose AgentGuard, which uses the agent orchestrator itself to
enumerate unsafe tool-use workflows and generate safety constraints that can
be validated before deployment~\cite{chen2025agentguard}.
These defences operate at the model fine-tuning or orchestration level;
PerceptFence operates at the input and output mediation layer without modifying
the underlying model, making it applicable to any assistant model but bounded
by the limits of rule-based detection, which the system design and evaluation
explicitly acknowledge.

%%%──────────────────────────────────────────────────────────────────────
\section{Design}\label{sec:design}
%%%──────────────────────────────────────────────────────────────────────

The runtime mediation layer is composed of seven modules along three control
surfaces:

\begin{itemize}
  \item \textbf{Observe} --- what the assistant model is permitted to see
        (screen capture adapter, speech event adapter, redaction engine,
        consent\,/\,policy engine).
  \item \textbf{Retain} --- what the session memory persists across turns
        (session memory gate).
  \item \textbf{Say} --- what the assistant is permitted to output
        (output guard).
\end{itemize}

The audit logger sits across all three surfaces and records decisions without
recording sensitive payload.

\subsection{Control surface 1 --- Observe}\label{sec:observe}

The observe surface controls what content the assistant may see. It is enforced
by two modules in sequence: the consent\,/\,policy engine and the redaction
engine.

\paragraph{Screen capture adapter.}
Captures display frames and extracts textual content. Output is untrusted text
plus structural metadata (window class, region, font size, recency). Does
\textbf{not} classify content; that is the consent\,/\,policy engine's job.

\paragraph{Speech event adapter.}
Captures audio events and emits utterance fragments with metadata (speaker
role, modality, timestamp). Consent decisions can be per-session,
per-modality, per-speaker, or per-content category; we adopt \textbf{per-modality
with per-content-category overrides}.

\paragraph{Consent\,/\,policy engine.}
Resolves the active consent profile per frame and per utterance. Inputs: window
class, modality, speaker role, content categories detected by the redactor, and
the user's session policy. Output: a policy action $\in$
\{\texttt{redact\_before\_model}, \texttt{suppress\_notification},
\texttt{summarize\_without\_identifier}, \texttt{block\_memory\_write},
\texttt{ignore\_screen\_instruction}, \texttt{require\_stable\_window},
\texttt{increase\_ocr\_sensitivity}, \texttt{selective\_redact}\}, plus the
rationale. The policy engine is the \textbf{only} authoritative decision-maker
for all three control surfaces. Policy downgrade is prevented by a
session-startup policy lock: relaxations require an authenticated re-consent
step, not a runtime toggle.

\paragraph{Redaction engine.}
Pattern-based and category-based redaction over extracted text and speech
fragments, applied after the policy decision. Replaces matched units with
stable placeholders (\texttt{[REDACTED]}, \texttt{[EMAIL]}, \texttt{[RECORD]},
\texttt{[PERSON]}). The engine is not model-based in this prototype, keeping
trust boundary TB2 deterministic; the design enforces \textbf{reduction of
sensitive exposure, not elimination}. Six deterministic transform families are
configured: (1)~Unicode-NFKC normalisation and confusable-character remapping
for homoglyph evasion; (2)~ignore-whitespace digit-group matching for split or
spaced PII; (3)~credential-pattern redaction for terminal tokens and
API-key-like strings; (4)~identifier summarisation with category placeholders;
(5)~prompt-injection neutralisation classifying screen-visible instruction text
as untrusted; and (6)~selective region redaction for mixed-sensitivity frames.

\subsection{Control surface 2 --- Retain}\label{sec:retain}

The retain surface controls what the assistant may remember across turns. It is
enforced by the session memory gate.

\paragraph{Session memory gate.}
Filters what the assistant's session memory persists. Default: no sensitive
units written to memory. When the policy action is
\texttt{block\_memory\_write}, the gate uses \textit{mode-A context exclusion}:
it zeros the model context for the affected turn before model invocation,
preventing sensitive content from reaching the model at all for that turn.
This is structurally stronger than conversation-history pruning, which removes
content from future turns but cannot prevent the model using content already in
the current context window. The cost is that benign task-critical units in the
excluded context are also unavailable for that turn (measured FBR cost in
Section~\ref{sec:per-fixture}). The gate records every exclusion event in the
audit log without storing the suppressed content.

\subsection{Control surface 3 --- Say}\label{sec:say}

The say surface controls what the assistant may output. It is enforced by the
output guard.

\paragraph{Output guard.}
Filters the assistant's response stream before display. Enforces output-side
policy actions (block, redact, paraphrase). Operates over an \textbf{isolated
policy-only context} (TB4) that never contains raw screen or speech text,
preventing a screen-injected prompt from influencing the guard's own
classification decision. Applies two filter stages: a literal-fragment denylist
for the current session's non-shareable category set, and an
indirect-disclosure classifier that detects outputs acknowledging the existence,
category, or format of gated content without verbatim repetition. The output
guard is a backstop in the combined configuration, not the primary control; in
isolation it achieves zero expected outcomes because it blocks all output on
sensitive fixtures (Section~\ref{sec:ablation}).

\subsection{Audit logger}\label{sec:audit}

Append-only log of policy decisions and module actions. Logs decision metadata
only---never raw screen text, raw speech, raw notification text, or sensitive
unit values. Each entry carries a SHA-256 hash chain for crash-evident integrity
(defence \textbf{not} claimed against A6).

\subsection{Implementation}\label{sec:impl}

The reference implementation is a standard-library-only Python package
(\texttt{screenshare\_mediator}) requiring Python $\geq$~3.10. No
machine-learning dependencies, network access, external APIs, or third-party
packages are required. All transforms are deterministic given the same fixture
and policy configuration, making runs bit-exact reproducible from the synthetic
fixture set.

\paragraph{Module inventory.}
Table~\ref{tab:modules} lists the eight Python modules.

\begin{table}[ht]
\caption{Module inventory of the reference implementation.}
\label{tab:modules}
\begin{tabular}{@{}lp{8.8cm}@{}}
\toprule
Module & Role \\
\midrule
\texttt{models.py}
  & Typed dataclasses (\texttt{CapturedSession}, \texttt{PolicyDecision},
    \texttt{MediatedContext}, \texttt{AuditEvent}, \texttt{GuardedOutput})
    shared across module boundaries. \\[3pt]
\texttt{capture.py}
  & Synthetic capture adapter; normalises a JSON fixture into a
    \texttt{CapturedSession} (TB1). \\[3pt]
\texttt{policy.py}
  & Consent\,/\,policy engine; maps \texttt{CapturedSession} to
    \texttt{PolicyDecision} using
    \texttt{consent\_redaction\_policy.json}. \\[3pt]
\texttt{redaction.py}
  & Six-family deterministic redaction engine; transforms raw context at
    TB2. \\[3pt]
\texttt{memory.py}
  & Mode-A session memory gate; enforces context exclusion and
    memory-write suppression at TB3. \\[3pt]
\texttt{output\_guard.py}
  & Rule-based output guard with isolated policy-only context (TB4). \\[3pt]
\texttt{audit.py}
  & Append-only SHA-256-chained audit logger (TA7). \\[3pt]
\texttt{runtime.py}
  & Baseline and guarded path composition; the two-path comparison
    interface used by the benchmark. \\
\bottomrule
\end{tabular}
\end{table}

\paragraph{Policy configuration.}
Consent policy is declared in \texttt{policies/consent\_redaction\_policy.json},
which maps scenario-class patterns to allowed action sets. The policy file is
the single configuration point; the policy engine reads it at startup and
applies it without code changes.

\paragraph{Synthetic fixture set.}
The fixture set covers eleven scenario classes. Each fixture is a JSON object
containing: scenario identifier, scenario class, synthetic screen text,
synthetic speech text, synthetic notification text (when applicable),
window-event metadata, annotated sensitive units ($U_i$), annotated benign
task-critical units ($Q_i$), expected policy action ($A_i$), and binary
task-success rubric answer ($y_i$). No fixture contains real screen captures,
real personal data, or real credentials.

\paragraph{Test suite.}
The test suite contains 29 tests: fixture-set coverage (6 tests), per-module
behaviour including all four adversarial-evasion paths (14 tests), benchmark
metric helpers (3 tests), audit hash-chain verification (4 tests), and a TB3
trust-boundary regression confirming context-excluded content does not appear in
the model context for the exclusion turn (1 test). All 29 pass on Python~3.14
and on Python~3.10.

\subsection{Design limits}\label{sec:designlimits}

\begin{itemize}
  \item Does not defend against A6, A7, or A8 (compromised infrastructure,
        supply chain, OS).
  \item Does not provide formal privacy guarantees.
  \item Does not perform model-based redaction; extended category coverage
        requires a redactor extension.
  \item Does not generalise beyond the seven modules without explicit
        threat-model re-review.
\end{itemize}

%%%──────────────────────────────────────────────────────────────────────
\section{Evaluation}\label{sec:eval}
%%%──────────────────────────────────────────────────────────────────────

\subsection{Metrics}\label{sec:metrics}

We define four primary metrics over a synthetic fixture set
$\mathcal{D} = \{x_1, \ldots, x_n\}$ with $n = 11$.
For each fixture, $U_i$ is the set of annotated synthetic sensitive units that
must not appear on prohibited surfaces (model context, assistant output, retained
memory, audit content beyond decision metadata), and $Q_i$ is the set of
annotated benign task-critical units that should remain available for task
completion.

\paragraph{Sensitive exposure rate (SER).}
For path $p \in \{B,\,G\}$ (baseline\,/\,guarded):
\begin{equation}
  \text{SER}(p) \;=\;
    \frac{\displaystyle\sum_{i}\sum_{u \in U_i} \text{exposed}(u,\,x_i,\,p)}
         {\displaystyle\sum_{i} |U_i|}
\end{equation}
where $\text{exposed}(u,\,x_i,\,p) = 1$ if a normalised match for $u$ appears
on a prohibited surface in path~$p$. Lower $\text{SER}(G)$ is better. Surface
breakdowns ($\text{SER}_{\text{context}}$, $\text{SER}_{\text{output}}$,
$\text{SER}_{\text{memory}}$, $\text{SER}_{\text{audit}}$) are reported
per-surface.

\paragraph{False block rate (FBR).}
Fraction of benign task-critical units $q \in Q_i$ that the guarded path blocks
or replaces in a way that prevents task completion.

\paragraph{Latency overhead (LOH).}
Wall-clock guarded-path runtime minus baseline runtime, on the same fixture,
under the same harness.

\paragraph{Task completion (TC).}
Rubric-based pass\,/\,fail per fixture, evaluated on assistant output.

\subsection{Benchmark composition}\label{sec:benchmark}

Eleven synthetic fixtures spanning seven non-adversarial classes (terminal
secret, chat notification, browser PII, spoken sensitive fragment, prompt
injection on screen, fast window switching, small-font zoomed UI) and three
adversarial-evasion classes (Unicode-homoglyph credential, split PII,
mixed-sensitivity scene). All fixtures are invented; no real screen captures,
real personal data, or production telemetry. Construction details and provenance
are in \texttt{src/data/synthetic/README.md} and
\texttt{supplement/artifact\_checklist.md}.

\subsection{Per-module ablation}\label{sec:ablation}

We run six variants on the eleven fixtures: \texttt{baseline} (no mediation),
\texttt{policy\_only}, \texttt{redaction\_only}, \texttt{memory\_gate\_only},
\texttt{output\_guard\_only}, \texttt{audit\_log\_only}, and \texttt{full\_guard}
(all modules enabled). Results from
\texttt{eval/results/per\_module\_ablation.csv} are shown in
Table~\ref{tab:ablation}.

\begin{table}[ht]
\caption{Per-module ablation results on eleven synthetic fixtures.
``Exp.\ rate'' = expected-outcome rate (fraction of fixtures reaching the
annotated expected outcome).}
\label{tab:ablation}
\resizebox{\textwidth}{!}{%
\begin{tabular}{@{}llrrrrrrr@{}}
\toprule
Variant & Modules & Exp.\ rate & \multicolumn{2}{c}{Exposures} & Out.\ & Audit & Mem.\ \\
        &         &            & Ctx. & Out. & blocks & events & writes \\
\midrule
\texttt{baseline}
  & none            & \textbf{0.000}~(0/11) & 10 & 10 &  0 &  0 &  0 \\
\texttt{policy\_only}
  & policy          & 0.000~(0/11)          & 10 &  8 &  0 &  0 &  0 \\
\texttt{redaction\_only}
  & policy+redact   & 0.909~(10/11)         &  1 &  1 &  0 &  0 &  0 \\
\texttt{memory\_gate\_only}
  & policy+memory   & 0.091~(1/11)          &  9 &  7 &  0 &  0 & 10 \\
\texttt{output\_guard\_only}
  & policy+guard    & 0.000~(0/11)          & 10 &  0 & 10 &  0 &  0 \\
\texttt{audit\_log\_only}
  & policy+audit    & 0.000~(0/11)          & 10 &  8 &  0 & 11 &  0 \\
\texttt{full\_guard}
  & all             & \textbf{1.000~(11/11)}&  0 &  0 &  5 & 23 & 10 \\
\bottomrule
\end{tabular}}
\end{table}

\paragraph{Reading the table.}
Redaction alone removes ten of eleven sensitive units from both context and
output; the missed fixture is \texttt{spoken\_sensitive\_fragment}, whose
expected policy action is \texttt{block\_memory\_write}---a retention control
that the redactor does not implement. Output guard alone blocks all assistant
outputs that would have leaked, but does so by blocking \textit{all} responses
on the affected fixtures---acceptable as a last-resort fallback, but not a
substitute for upstream filtering. Memory gate alone prevents persistence but
lets the model see and respond with sensitive content. The full guard
composition reaches the expected outcome on every fixture in the benchmark,
with no context exposures, no output exposures, five precise output blocks
(used only when upstream redaction was insufficient), an audit trail of 23
decision events, and ten memory-gate writes recorded by the audit logger.

\subsection{Per-fixture results and adversarial-evasion
            analysis}\label{sec:per-fixture}

The three adversarial-evasion fixtures are \texttt{homoglyph\_credential},
\texttt{split\_pii}, and \texttt{encoded\_screen\_instruction}. Results from
\texttt{eval/results/per\_fixture\_ablation.csv} are summarised in
Table~\ref{tab:evasion}.

\begin{table}[ht]
\caption{Adversarial-evasion fixture outcomes.}
\label{tab:evasion}
\begin{tabular}{@{}p{2.8cm}llll@{}}
\toprule
Fixture & redact.-only & full-guard & Out.\ blk & Primary mechanism \\
\midrule
\texttt{homoglyph\_} \texttt{credential}
  & met & met & No
  & Unicode normalisation before pattern check removes the homoglyph
    substitution; output guard not invoked. \\[4pt]
\texttt{split\_pii}
  & met & met & Yes
  & Whitespace\,/\,separator normalisation catches split digits;
    output guard fires as defence-in-depth. \\[4pt]
\texttt{encoded\_screen\_} \texttt{instruction}
  & met & met & Yes
  & Redactor strips encoding marker; output guard additionally blocks
    as defence-in-depth. \\
\bottomrule
\end{tabular}
\end{table}

All three adversarial-evasion fixtures are handled by \texttt{full\_guard}.
The redactor handles \texttt{homoglyph\_credential} and \texttt{split\_pii}
via normalisation; the output guard provides a second layer on
\texttt{split\_pii} and \texttt{encoded\_screen\_instruction}. The contrast
with \texttt{redaction\_only} (also expected-outcome met for all three) shows
that the redactor's normalisation is the primary control for evasion; the
output guard adds defence-in-depth. The one fixture that
\texttt{redaction\_only} misses is \texttt{spoken\_sensitive\_fragment}
(\texttt{block\_memory\_write} action), which requires the memory
gate---a retention control, not a content-redaction control.

\subsection{Latency overhead}\label{sec:latency}

Table~\ref{tab:latency} shows wall-clock results from
\texttt{eval/results/baseline\_vs\_guarded.csv} (200 paired iterations per
fixture).

\begin{table}[ht]
\caption{Latency overhead (200 paired iterations per fixture; values in
         microseconds, $\mu$s).}
\label{tab:latency}
\begin{tabular}{@{}lrrrr@{}}
\toprule
Fixture & Baseline med. & Guarded med. & Guarded p99 & Overhead ratio \\
\midrule
\texttt{terminal\_secret}             & 2.2 & 32.8 &  40.0 & 13.9$\times$ \\
\texttt{chat\_notification}           & 2.2 & 55.5 &  73.1 & 24.6$\times$ \\
\texttt{browser\_pii}                 & 2.1 & 35.4 &  46.0 & 15.7$\times$ \\
\texttt{spoken\_sensitive\_fragment}  & 2.0 & 35.9 &  85.9 & 17.0$\times$ \\
\texttt{prompt\_injection\_on\_screen}& 2.1 & 32.6 &  44.8 & 14.6$\times$ \\
\texttt{fast\_window\_switching}      & 2.4 & 32.8 &  43.3 & 12.8$\times$ \\
\texttt{small\_font\_zoomed\_ui}      & 2.1 & 28.8 &  99.8 & 12.5$\times$ \\
\texttt{homoglyph\_credential}        & 2.1 & 27.3 &  28.1 & 11.9$\times$ \\
\texttt{split\_pii}                   & 2.1 & 28.6 &  37.1 & 12.5$\times$ \\
\texttt{mixed\_sensitivity}           & 2.1 & 56.8 &  81.3 & 25.5$\times$ \\
\texttt{encoded\_screen\_instruction} & 2.1 & 71.1 & 103.6 & 32.3$\times$ \\
\bottomrule
\end{tabular}
\end{table}

The median guarded latency across all fixtures is approximately 32.8\,$\mu$s
(range 27.3--71.1\,$\mu$s). The p99 reaches 103.6\,$\mu$s for
\texttt{encoded\_screen\_instruction}, which requires the most complex
output-guard evaluation. The baseline median is approximately 2.2\,$\mu$s. The
overhead ratio ranges from 11.9$\times$ to 32.3$\times$, reflecting that the
mediation layer dominates harness time relative to the no-op baseline, but the
absolute overhead remains sub-millisecond in every case.

These figures reflect a single-process Python evaluation harness on synthetic
fixtures. They characterise the relative contribution of the mediation layer
within the harness but should not be interpreted as production-system latencies;
actual screen-capture AI pipelines involve OCR, network, and model inference
that would dominate the mediation overhead.

\subsection{Threats to validity}\label{sec:threats}

\begin{itemize}
  \item \textbf{Synthetic-only.} All fixtures are invented. Production-scale
        generalisation is not claimed.
  \item \textbf{Eleven fixtures.} Coverage is intentionally bounded;
        adversarial-evasion classes beyond the three studied are not
        characterised.
  \item \textbf{English-only synthetic content.} Multi-lingual evasion is not
        studied.
  \item \textbf{Single-system evaluation.} No comparison to prior
        runtime-mediation systems because, to our knowledge, none with the same
        control-surface decomposition exists. The closest prior systems operate
        at the OS permission layer (TaintDroid, Android permission model) or
        session-level controls (Microsoft Recall, Zoom AI Companion), not at the
        per-frame, per-content-category runtime level described here.
\end{itemize}

%%%──────────────────────────────────────────────────────────────────────
\section{Discussion}\label{sec:discussion}
%%%──────────────────────────────────────────────────────────────────────

The ablation supports a small but specific design claim: \textbf{no single
module is sufficient, and the four mediation modules (redaction, memory gate,
output guard, audit logger) compose under the policy engine to achieve full
expected-outcome on the synthetic benchmark.}
The result is consistent with a defence-in-depth architecture: redaction handles
the common path, output guard catches what redaction misses, memory gate prevents
cross-turn leakage, and audit logger preserves decision provenance.

\paragraph{Deployment considerations.}
The seven-module decomposition assumes a controlled assistant runtime where the
developer controls capture, redaction, model invocation, memory, and output.
SaaS assistants without this control cannot adopt the design as-is. Hybrid
deployments (browser extension + assistant API) would require redaction at the
extension boundary and trust-model adjustments.

\paragraph{What we did not measure.}
Real-user task-completion rates, latency on production captures, additional
adversarial classes beyond the three studied, multi-lingual performance.

%%%──────────────────────────────────────────────────────────────────────
\section{Limitations}\label{sec:limitations}
%%%──────────────────────────────────────────────────────────────────────

\begin{itemize}
  \item Synthetic benchmark only; no real-user evaluation.
  \item Eleven fixtures; adversarial classes beyond the three studied are
        uncharacterised.
  \item Pattern-based redactor; additional content categories require a
        redactor extension.
  \item A6, A7, and A8 explicitly excluded from the threat model.
  \item No formal privacy guarantee.
\end{itemize}

%%%──────────────────────────────────────────────────────────────────────
\section{Ethics}\label{sec:ethics}
%%%──────────────────────────────────────────────────────────────────────

All evaluation data is synthetic. No real screen captures, real personal data,
real audio recordings, or real notification content was collected, used, or
stored. No human subjects were involved. IRB approval is not applicable. The
synthetic fixture set is documented under \texttt{src/data/synthetic/README.md}.

The design does not enable surveillance: it reduces, not increases, the
assistant's access to sensitive content. The design does not weaken existing
privacy controls: it composes with platform-level controls and operates
strictly inside the assistant runtime.

%%%──────────────────────────────────────────────────────────────────────
\section*{Declarations}
%%%──────────────────────────────────────────────────────────────────────

\paragraph{Funding.}
The author received no external funding for this work.

\paragraph{Competing interests.}
The author is employed by Parafin, Inc. The study uses only synthetic data and
does not use Parafin customer data, production telemetry, proprietary code, or
infrastructure. The author declares no other competing interests.

\paragraph{Ethics approval and consent to participate.}
Not applicable. The study used synthetic fixtures and synthetic evaluation
outputs only. It did not involve human participants, human-subject data, real
screen captures, real audio recordings, customer data, production telemetry, or
organization-specific datasets.

\paragraph{Consent for publication.}
Not applicable. The manuscript contains no identifiable human-subject data,
real screenshots, real notifications, real speech recordings, customer records,
or private organizational data.

\paragraph{Data availability.}
The synthetic fixture definitions and evaluation result CSV files supporting
the reported findings are included with the supplementary materials. No real
screen captures, personal data, customer data, production telemetry, or
human-subject data were used.

\paragraph{Code availability.}
The reference implementation can be supplied as supplementary material with the
manuscript, subject to the license terms included with the artifact package.
Public repository or archive availability should be asserted only after the
artifact publication and license posture are finalized.

\paragraph{Author contributions.}
Neeraj Kumar Singh is the sole author. He designed the runtime mediation
architecture, implemented the synthetic prototype and evaluation fixtures,
analyzed the results, and wrote the manuscript.

\paragraph{Generative AI / tool-use disclosure.}
AI-assisted tools were used for editorial drafting, formatting support, coding
support, and submission-package preparation. They were not credited as authors.
The author reviewed and remains responsible for the manuscript content, claims,
code, synthetic data posture, and final submission declarations.

%%%──────────────────────────────────────────────────────────────────────
\section{Conclusion}\label{sec:conclusion}
%%%──────────────────────────────────────────────────────────────────────

We presented a runtime mediation layer for live screen-share AI assistants,
decomposed into seven modules along three control surfaces (observe, retain,
say) and motivated by an explicit threat model.
A per-module ablation on a synthetic benchmark of eleven scenario classes shows
that the composed full guard achieves expected outcomes on every fixture while
the unmediated baseline reaches it on none, and that no single module is
individually sufficient.
We make bounded claims tied to a synthetic benchmark and explicitly exclude
formal privacy, host compromise, and supply-chain threats from our scope.

%%%──────────────────────────────────────────────────────────────────────
%%% Bibliography
%%%──────────────────────────────────────────────────────────────────────

\bibliographystyle{sn-mathphys-num}
\bibliography{../references}

\end{document}
